\section{Renormalization fixed points: a shared adapted basis}%
\label{sec:asymex_rfp_basis}

The two examples considered below use stacked products of matrix product
operator tensors whose bond spaces are tensor squares. Reordering the
stacked bond coordinates separates a four-dimensional inner bond from
the outer bond space. The inner bond consists of the two half-bonds
joined by summing over the intermediate physical index. Both examples
use the same adapted basis: two Bell vectors and the two coordinate
vectors orthogonal to them.

The length dependence comes from powers of the weights attached to
rank-one idempotents. In the notation of the RFP algebra, the structure
coefficients are
$c_{\alpha,\beta,\gamma}^{(L)}
=\tr(\chi_{\alpha,\beta,\gamma}^L)$.
Each weighted copy of a target contributes its weight to the power $L$;
weights are summed when they multiply the same target. This is the
power-sum mechanism of
Theorem~\ref{thm:asym_power_sum_coefficients}.
The results below establish the tensor decompositions and periodic trace
identities for these examples, rather than a classification of
length-dependent RFPs.

\begin{definition}[The idempotents of an adapted basis]
	\label{def:asymex_basis_proj}
	\lean{MPSTensor.basisProj, P6Compression.pairReorder, P6Compression.pairU,
		P6Compression.pairUinv, P6Compression.pairProjInt,
		P6Compression.pairBlockInt}
	\leanok
	\uses{prop:asym_projector_weighted_sum}
	For mutually inverse matrices $U$ and $U^{-1}$, define
	$\Pi_j=U^{-1}E_{jj}U$. This is the product of the $j$-th column of
	$U^{-1}$ with the $j$-th row of $U$.
	In both examples below, a stacked bond coordinate is a quadruple of
	bits. Reorder it into the inner bond formed by the two middle bits
	and the outer bond formed by the two outer bits. Choose the adapted
	basis of the inner bond whose first two vectors are
	$(1,0,0,\pm1)$ and whose last two vectors are the remaining coordinate
	vectors. The two idempotents associated with the Bell vectors carry
	the weighted target tensors. Tensor products of these idempotents
	with the target matrices are then expressed in the original stacked
	bond coordinates.
\end{definition}

\begin{lemma}[The idempotents of an adapted basis are mutually orthogonal]
	\label{lem:asymex_basis_proj}
	\lean{MPSTensor.basisProj_mul_basisProj_self,
		MPSTensor.basisProj_mul_basisProj_of_ne, MPSTensor.mul_basisProj_mul}
	\leanok
	\uses{def:asymex_basis_proj}
	The idempotents satisfy $\Pi_j^2=\Pi_j$,
	$\Pi_j\Pi_k=0$ for $j\ne k$, and $U\Pi_jU^{-1}=E_{jj}$.
\end{lemma}

\begin{proof}
	\leanok
	The idempotent is $U^{-1}E_{jj}U$, so conjugation by $U$ returns
	$E_{jj}$, and the products reduce to $E_{jj}E_{kk}$, which is $E_{jj}$
	for $j=k$ and zero otherwise.
\end{proof}

\begin{theorem}[From an integer matrix identity to a weighted direct sum]
	\label{thm:asymex_stacked_pair}
	\lean{P6Compression.stackedPair_letter_identity}
	\leanok
	\uses{def:asymex_basis_proj, prop:asym_projector_weighted_sum}
	Suppose the stacked tensor is $c_B$ times an integer tensor
	$\widehat B$, and each target tensor $C_s$ is $c_C$ times an integer
	tensor $\widehat C_s$, with $s\in\{0,1\}$. Suppose there are a nonzero
	integer $n$ and integers $k_s$ such that, for every $s$ and every
	physical index $a$,
	\begin{align}
		n\,\mu_s\,\frac{c_C}{2} & =k_s\,c_B,
		\label{eq:asymex_stacked_pair_scalar}                            \\
		n\,\widehat B^a         & =\sum_sk_s\,\bigl(\widehat\Pi_s\otimes
		\widehat C_s^a\bigr),
		\label{eq:asymex_stacked_pair_letters}
	\end{align}
	where $\widehat\Pi_s=2\Pi_s$ and the right-hand side of
	(\ref{eq:asymex_stacked_pair_letters}) is expressed in the stacked
	bond coordinates. Then
	$B^a=\sum_s\mu_s(\Pi_s\otimes C_s^a)$ in those coordinates, which is
	the hypothesis of Proposition~\ref{prop:asym_projector_weighted_sum}.
\end{theorem}

\begin{proof}
	\leanok
	\uses{lem:asymex_basis_proj}
	Expand the right-hand side entrywise using the idempotents of
	Lemma~\ref{lem:asymex_basis_proj}. Since
	$\Pi_s=\widehat\Pi_s/2$ and $C_s^a=c_C\widehat C_s^a$,
	its $s$-th summand is
	$\mu_sc_C(\widehat\Pi_s\otimes\widehat C_s^a)/2$.
	Multiplying by $n$ and applying
	(\ref{eq:asymex_stacked_pair_scalar}) and
	(\ref{eq:asymex_stacked_pair_letters}) gives
	$c_Bn\widehat B^a=nB^a$. Cancel the nonzero scalar $n$.
\end{proof}

\begin{theorem}[Normality from scaled matrix units]
	\label{thm:asymex_normal_units}
	\lean{P6Compression.isNormal_of_single_eq_smul}
	\leanok
	\uses{def:normal}
	Let a tensor be one nonzero scalar times an integer tensor, and suppose
	that for every position $(x,y)$ of its bond algebra some matrix of the
	integer tensor is a nonzero integer multiple of $E_{xy}$. Then the
	one-site matrices of the original tensor span the full matrix algebra
	of the bond space, so the tensor is normal at blocking length one in
	the sense of Definition~\ref{def:normal}.
	This conclusion does not impose spectral-radius normalization;
	the term NT in the RFP canonical form additionally requires that
	normalization.
\end{theorem}

\begin{proof}
	\leanok
	Dividing the designated letter by its nonzero scalar factor exhibits
	each matrix unit in the span of the one-site matrices, and the matrix
	units span the whole algebra.
\end{proof}

\section{The one-label fixed point}%
\label{sec:asymex_one_label}

The first example has four one-site matrices:
\begin{align}
	A^0 & =\tfrac45E_{00},
	\label{eq:asymex_one_label_matrices} \\
	A^1 & =\tfrac45E_{11}, \\
	A^2 & =\tfrac35E_{01}, \\
	A^3 & =\tfrac35E_{10}.
\end{align}
The squared scales are $(1\pm\lambda)/2$ for $\lambda=7/25$.
Reading the tensor in the vertical direction gives
$M^{(i,j)}=A^i\otimes A^j$, with sixteen physical index pairs.
Its stacked product factorizes as $A^i\otimes T\otimes A^k$, where
$T=\sum_jA^j\otimes A^j$ acts on the inner bond and has spectrum
$\{1,7/25,0,0\}$. Here $T$ denotes this matrix, not a tpCPM in the
definition of an RFP. Its spectral decomposition gives two weighted
copies of the same target tensor, with weights $1$ and $7/25$.
Thus the associated periodic structure coefficient is
$c^{(L)}=\tr(\chi^L)=1+(7/25)^L$ with
$\chi=\diag(1,7/25)$, rather than a sum of contributions from
inequivalent target tensors.

\begin{definition}[The one-label tensors]
	\label{def:asymex_one_label}
	\lean{P6Compression.oneLabelM, P6Compression.oneLabelTarget,
		P6Compression.oneLabelStacked, P6Compression.oneLabelBlocks,
		P6Compression.oneLabelWeights}
	\leanok
	\uses{def:mpo_tensor, def:mpo_to_mps, def:mpdo_operator_product}
	The tensor read in the vertical direction is
	$M^{(i,j)}=A^i\otimes A^j$, with the matrices
	(\ref{eq:asymex_one_label_matrices}), of bond dimension four.
	Viewed as a matrix product tensor over the sixteen physical index
	pairs, it is the target tensor. Let $B$ be its stacked product with
	itself, of bond dimension sixteen. The two target tensors are copies
	of $M$ with weights $\mu_0=1$ and $\mu_1=7/25$.
\end{definition}

\begin{center}
	% Formula: the periodic contraction of the vertical tensor $M$ with
	% the inner-bond state $\sigma=\sum_b\lambda_b\Pi_b$ inserted on
	% every bond, whose weights are the spectrum $\{1,7/25,0,0\}$ of
	% $T=\sum_jA^j\otimes A^j$ (\S\ref{sec:asymex_rfp_basis}); this is
	% the spectral mechanism behind
	% Theorem~\ref{thm:asymex_one_label_compression}.
	% Source: \S\ref{sec:asymex_one_label}, introductory paragraph.
	% Ink-to-index map: the two explicit dot nodes are the vertical
	% tensor $M$, each carrying its physical leg of the pair alphabet;
	% each ring is the inner-bond state $\sigma$, with no physical leg
	% of its own; the trailing dots abbreviate the remaining alternating
	% $M$-$\sigma$ pairs of the periodic ring.
	% Contraction set: every virtual bond, alternating $M$ and $\sigma$,
	% closed around the ring; no leg is left open.
	% Boundary signature: (virtual-west=0 | virtual-east=0 |
	% physical-up=2 | physical-down=0) for the finite schematic shown.
	\tenkzkernel
	\begin{tenkz}[cols=6, boundary=periodic, physical=up]
		\tn[label pos=135, ports={90:physical:$a_1$}]{M} &
		\tn[skin=ring]{\sigma} &
		\tn[label pos=135, ports={90:physical:$a_2$}]{M} &
		\tn[skin=ring]{\sigma} &
		\tn[skin=dots]{} & \tn[skin=ring]{\sigma}
	\end{tenkz}
	\qquad $\sigma=\sum_b\lambda_b\,\Pi_b$
\end{center}

\begin{theorem}[Normality of the one-label target]
	\label{thm:asymex_one_label_normal}
	\lean{P6Compression.oneLabelTarget_isNormal}
	\leanok
	\uses{def:asymex_one_label, def:normal}
	The vertical tensor is normal at blocking length one in the sense of
	Definition~\ref{def:normal}: its sixteen matrices are scaled matrix
	units covering every position of the four-by-four matrix algebra.
\end{theorem}

\begin{proof}
	\leanok
	\uses{thm:asymex_normal_units}
	Each of the sixteen positions is realized by one of the sixteen letters
	with a nonzero integer factor, a finite verification over integer
	matrices; Theorem~\ref{thm:asymex_normal_units} then gives normality.
\end{proof}

\begin{theorem}[Split compression of the one-label fixed point]
	\label{thm:asymex_one_label_compression}
	% Principal declaration: P6Compression.oneLabel_isReduction.
	\lean{P6Compression.oneLabelCompression, P6Compression.oneLabel_isReduction,
		P6Compression.oneLabel_left_mul_right_of_ne, P6Compression.oneLabel_z_eq,
		P6Compression.oneLabel_dim_eq, P6Compression.oneLabel_remainder,
		P6Compression.oneLabel_mul_right_eq_right_mul,
		P6Compression.oneLabel_left_mul_eq_mul_left}
	\leanok
	\uses{def:asymex_one_label, thm:asym_multiblock}
	The stacked tensor $B$ of Definition~\ref{def:asymex_one_label} has the
	block triangular form of Theorem~\ref{thm:asym_multiblock} for the two
	weighted copies $C_s=\mu_sM$, with an eight-dimensional zero
	complement, $z=8$, and $16=4+4+8$. The compression pairs are mutually
	biorthogonal and satisfy $W_sB^wV_s=C_s^w$ for every word.
	The remainder vanishes identically, and the compression maps
	intertwine the one-site matrices in both directions:
	$B^aV_s=V_s(\mu_sM^a)$ and $W_sB^a=(\mu_sM^a)W_s$.
\end{theorem}

\begin{proof}
	\leanok
	\uses{thm:asymex_stacked_pair, prop:asym_projector_weighted_sum}
	Multiplying the stacked tensor by two and expanding it in the adapted
	basis of Section~\ref{sec:asymex_rfp_basis} gives the integer matrix
	identity of Theorem~\ref{thm:asymex_stacked_pair} with the integer
	coefficients $25$ and $7$, a finite verification over the sixteen
	physical indices. The matching scalar identity holds because the two
	weights are $1$ and $7/25$.
	Proposition~\ref{prop:asym_projector_weighted_sum} then supplies the
	compression maps, the vanishing remainder and the one-site
	intertwining identities, with a zero complement of dimension
	$8=2\cdot4$.
\end{proof}

\begin{center}
	% Formula: B^aV_s=V_s(\mu_sM^a) and W_sB^a=(\mu_sM^a)W_s of
	% Theorem~\ref{thm:asymex_one_label_compression}, for the weights
	% $\mu_s\in\{1,7/25\}$.
	% Source: Theorem~\ref{thm:asymex_one_label_compression}.
	% Ink-to-index map: the ring is the compression map $V_s$ or $W_s$;
	% the square node is the stacked source $B$ or the weighted target
	% $\mu_sM$.
	% Contraction set: the single virtual bond joining the ring to the
	% square node in each pair; the physical pair leg $a$ remains open.
	% Boundary signature: (virtual-west=1 | virtual-east=1 |
	% physical-up=1 | physical-down=0) for each pictured equality.
	\tenkzkernel
	\begin{tenkz}[rows={wire}, cols=2, west=open, east=open]
		\tn[label pos=s, ports={90:physical:$a$}]{B} & \tn[skin=ring]{V_s}
	\end{tenkz}
	\quad $=$ \quad
	\begin{tenkz}[rows={wire}, cols=2, west=open, east=open]
		\tn[skin=ring]{V_s} & \tn[label pos=s, ports={90:physical:$a$}]{\mu_sM}
	\end{tenkz}
	\\[2ex]
	\begin{tenkz}[rows={wire}, cols=2, west=open, east=open]
		\tn[skin=ring]{W_s} & \tn[label pos=s, ports={90:physical:$a$}]{B}
	\end{tenkz}
	\quad $=$ \quad
	\begin{tenkz}[rows={wire}, cols=2, west=open, east=open]
		\tn[label pos=s, ports={90:physical:$a$}]{\mu_sM} & \tn[skin=ring]{W_s}
	\end{tenkz}
\end{center}

\begin{theorem}[The one-label structure coefficient]
	\label{thm:asymex_one_label_trace}
	\lean{P6Compression.oneLabel_trace_evalWord}
	\leanok
	\uses{def:asymex_one_label}
	For every nonempty word $w$ over the sixteen-letter alphabet,
	\begin{align}
		\tr(B^w) & =\left(1+(\tfrac{7}{25})^{|w|}\right)\tr(M^w).
		\label{eq:asymex_one_label_trace}
	\end{align}
\end{theorem}

\begin{proof}
	\leanok
	\uses{thm:asymex_one_label_compression, prop:asym_compression_trace}
	Proposition~\ref{prop:asym_compression_trace} applied to the compression
	in Theorem~\ref{thm:asymex_one_label_compression} gives the sum of the
	two weighted target traces. Pulling each weight out of its word
	evaluation contributes $1$ and $(7/25)^{|w|}$.
\end{proof}

\section{The graded quantum-dimer twist}%
\label{sec:asymex_dimer}

In the two-sector graded quantum-dimer twist, a horizontal bond index is a
triple $(p,p',k)$ of bits: the ket and bra indices of the left half-bond
and a sector label. The vertical site space is $\C^2\otimes\C^2$.
Each nonzero matrix of the sector-$f$ tensor is a scaled matrix unit
with scale $\tfrac12C_k[p][p']\tau_k[f]$, where
$C_0=xP_++yP_-$, $C_1=xP_+-yP_-$,
$\tau_0=(1,1)$ and $\tau_1=(1,-1)$.
Matrices whose two horizontal bond indices have different sector labels
vanish.

The target tensors for the stacked product of sectors $f$ and $f'$ are
the sectors $f+f'$ and $f+f'+1$, with addition modulo two, and their
weights are $x/2$ and $y/2$. At $x=7/8$, $y=1/8$, these weights are
$7/16$ and $1/16$. Unlike the one-label example, the two contributions
multiply different target tensors. For the product of sector zero with
itself, the periodic structure coefficients are
$c_{0,0,0}^{(L)}=(7/16)^L$ and
$c_{0,0,1}^{(L)}=(1/16)^L$, the power sums of the one-dimensional
positive diagonal matrices $\chi_{0,0,0}=(7/16)$ and
$\chi_{0,0,1}=(1/16)$, respectively.
The compression and trace theorems below concern this pair of sectors.

\begin{definition}[The sector tensors and their stacked product]
	\label{def:asymex_dimer}
	\lean{P6Compression.dimerM, P6Compression.dimerTarget,
		P6Compression.dimerStacked, P6Compression.dimerBlocks,
		P6Compression.dimerWeights}
	\leanok
	\uses{def:mpo_tensor, def:mpo_to_mps, def:mpdo_operator_product}
	At $x=7/8$ and $y=1/8$, the sector-$f$ tensor has, at the index pair
	$(a,b)$ of horizontal bonds, the scaled matrix unit whose row is the
	pair of ket bits of $a$ and $b$, whose column is the pair of bra bits,
	and whose scale is $\tfrac12C_k[p][p']\tau_k[f]$ for their common
	sector label $k$. It vanishes when the two sector labels differ.
	Read in the vertical direction, over the sixty-four index pairs,
	this gives the target tensor $M_f$.
	The stacked product of the sector-$f$ and sector-$f'$ tensors has
	bond dimension sixteen. The two target tensors are $M_{f+f'}$ and
	$M_{f+f'+1}$ with weights $7/16$ and $1/16$, respectively, where
	sector addition is modulo two.
\end{definition}

\begin{center}
	% Formula: the periodic contraction of the sector-$f$ tensor $M_f$
	% with the inner-bond state $\sigma=\sum_b\lambda_b\Pi_b$ inserted
	% on every bond, the same adapted-basis mechanism of
	% \S\ref{sec:asymex_rfp_basis} that gives the weights $7/16$ and
	% $1/16$.
	% Source: \S\ref{sec:asymex_dimer}, introductory paragraph.
	% Ink-to-index map: the two explicit dot nodes are the sector tensor
	% $M_f$, each carrying its physical leg of the pair alphabet; each
	% ring is the inner-bond state $\sigma$, with no physical leg of its
	% own; the trailing dots abbreviate the remaining alternating
	% $M_f$-$\sigma$ pairs of the periodic ring.
	% Contraction set: every virtual bond, alternating $M_f$ and
	% $\sigma$, closed around the ring; no leg is left open.
	% Boundary signature: (virtual-west=0 | virtual-east=0 |
	% physical-up=2 | physical-down=0) for the finite schematic shown.
	\tenkzkernel
	\begin{tenkz}[cols=6, boundary=periodic, physical=up]
		\tn[label pos=135, ports={90:physical:$a_1$}]{M_f} &
		\tn[skin=ring]{\sigma} &
		\tn[label pos=135, ports={90:physical:$a_2$}]{M_f} &
		\tn[skin=ring]{\sigma} &
		\tn[skin=dots]{} & \tn[skin=ring]{\sigma}
	\end{tenkz}
	\qquad $\sigma=\sum_b\lambda_b\,\Pi_b$
\end{center}

\begin{theorem}[Normality of the sectors]
	\label{thm:asymex_dimer_normal}
	\lean{P6Compression.dimerTarget_isNormal}
	\leanok
	\uses{def:asymex_dimer, def:normal}
	Each sector tensor is normal at blocking length one in the sense of
	Definition~\ref{def:normal}: its thirty-two nonzero matrices are
	scaled matrix units covering every position of the four-by-four
	matrix algebra.
\end{theorem}

\begin{proof}
	\leanok
	\uses{thm:asymex_normal_units}
	Sixteen letters lying in the sector label zero, where the scale does
	not depend on $f$, realize the sixteen positions with nonzero integer
	factors, a finite verification over integer matrices;
	Theorem~\ref{thm:asymex_normal_units} then gives normality.
\end{proof}

\begin{theorem}[Split compression of the dimer fusion]
	\label{thm:asymex_dimer_compression}
	% Principal declaration: P6Compression.dimer_isReduction.
	\lean{P6Compression.dimerCompression, P6Compression.dimer_isReduction,
		P6Compression.dimer_left_mul_right_of_ne, P6Compression.dimer_z_eq,
		P6Compression.dimer_dim_eq, P6Compression.dimer_remainder,
		P6Compression.dimer_mul_right_eq_right_mul,
		P6Compression.dimer_left_mul_eq_mul_left}
	\leanok
	\uses{def:asymex_dimer, thm:asym_multiblock}
	For the product of sector zero with itself, let $B$ be the stacked
	tensor of Definition~\ref{def:asymex_dimer}.
	It has the block triangular form of
	Theorem~\ref{thm:asym_multiblock} for the two weighted target tensors
	$C_s=\mu_sM_s$, where $\mu_0=7/16$ and $\mu_1=1/16$, with an
	eight-dimensional zero complement, $z=8$, and $16=4+4+8$.
	The compression pairs are mutually biorthogonal and satisfy
	$W_sB^wV_s=C_s^w$ for every word. The remainder vanishes identically,
	and the compression maps intertwine the one-site matrices in both
	directions.
\end{theorem}

\begin{proof}
	\leanok
	\uses{thm:asymex_stacked_pair, prop:asym_projector_weighted_sum}
	A matrix of the stacked tensor vanishes unless its two outer
	horizontal bond indices carry the same sector label. For a matching
	pair, the sum over the intermediate index runs over the four bond
	indices of that label. Multiplying the resulting matrix by two and
	expanding it in the adapted basis of
	Section~\ref{sec:asymex_rfp_basis} gives the integer matrix identity
	of Theorem~\ref{thm:asymex_stacked_pair} with coefficients $7$ and
	$1$, a finite verification over the sixty-four physical indices in
	each of the two sector-label cases. The matching scalar identity
	holds because the weights are $7/16$ and $1/16$.
	Proposition~\ref{prop:asym_projector_weighted_sum} then supplies the
	compression maps, the vanishing remainder and the one-site
	intertwining identities.
\end{proof}

\begin{center}
	% Formula: the sitewise intertwiners B^aV_s=V_s(\mu_sM_s^a) and
	% W_sB^a=(\mu_sM_s^a)W_s of Theorem~\ref{thm:asymex_dimer_compression},
	% for the two weighted target sectors $s\in\{f+f',f+f'+1\}$ with
	% weights $\mu_s\in\{7/16,1/16\}$.
	% Source: Theorem~\ref{thm:asymex_dimer_compression}.
	% Ink-to-index map: the ring is the compression map $V_s$ or $W_s$;
	% the square node is the stacked source $B$ or the weighted sector
	% target $\mu_sM_s$.
	% Contraction set: the single virtual bond joining the ring to the
	% square node in each pair; the physical pair leg $a$ remains open.
	% Boundary signature: (virtual-west=1 | virtual-east=1 |
	% physical-up=1 | physical-down=0) for each pictured equality.
	\tenkzkernel
	\begin{tenkz}[rows={wire}, cols=2, west=open, east=open]
		\tn[label pos=s, ports={90:physical:$a$}]{B} & \tn[skin=ring]{V_s}
	\end{tenkz}
	\quad $=$ \quad
	\begin{tenkz}[rows={wire}, cols=2, west=open, east=open]
		\tn[skin=ring]{V_s} &
		\tn[label pos=s, ports={90:physical:$a$}]{\mu_sM_s}
	\end{tenkz}
	\\[2ex]
	\begin{tenkz}[rows={wire}, cols=2, west=open, east=open]
		\tn[skin=ring]{W_s} & \tn[label pos=s, ports={90:physical:$a$}]{B}
	\end{tenkz}
	\quad $=$ \quad
	\begin{tenkz}[rows={wire}, cols=2, west=open, east=open]
		\tn[label pos=s, ports={90:physical:$a$}]{\mu_sM_s} & \tn[skin=ring]{W_s}
	\end{tenkz}
\end{center}

\begin{theorem}[The dimer structure coefficients]
	\label{thm:asymex_dimer_trace}
	\lean{P6Compression.dimer_trace_evalWord}
	\leanok
	\uses{def:asymex_dimer}
	For the product of sector zero with itself and every nonempty word
	$w$ over the sixty-four-letter alphabet,
	\begin{align}
		\tr(B^w)
		 & =(\tfrac{7}{16})^{|w|}\tr(M_0^w)
		+(\tfrac{1}{16})^{|w|}\tr(M_1^w),
		\label{eq:asymex_dimer_trace}
	\end{align}
	where $M_0$ and $M_1$ are the two sector tensors.
\end{theorem}

\begin{proof}
	\leanok
	\uses{thm:asymex_dimer_compression, prop:asym_compression_trace}
	Proposition~\ref{prop:asym_compression_trace} applied to the compression
	in Theorem~\ref{thm:asymex_dimer_compression} gives the sum of the two
	weighted target traces. Pulling each weight out of its word
	evaluation gives the factors $(7/16)^{|w|}$ and $(1/16)^{|w|}$ in
	(\ref{eq:asymex_dimer_trace}).
\end{proof}

\begin{remark}[Scope of the dimer fusion identities]
	\label{rem:asymex_dimer_scope}
	Theorems~\ref{thm:asymex_dimer_compression}
	and~\ref{thm:asymex_dimer_trace} concern the product of sector zero
	with itself. The remaining three pairs of sectors satisfy the same
	one-site matrix identity, with the same adapted basis and weights,
	and with targets $M_{f+f'}$ and $M_{f+f'+1}$. For each pair this
	requires a separate verification over the sixty-four physical
	indices, which is not carried out here.
\end{remark}
